% ============================================================
% Preamble: packages, layout and shared styles for the thesis
% ============================================================

\usepackage[utf8]{inputenc}
\usepackage[T1]{fontenc}
\usepackage{lmodern}
\usepackage[italian,english]{babel}   % English is the main document language;
                                       % Italian is only used for the cover page
\usepackage[a4paper,margin=2.5cm]{geometry}
\usepackage{pdflscape}   % landscape pages for full-page figures (e.g. the Qsys block diagram)

\usepackage{indentfirst}   % also indent the first paragraph after a heading
\setlength{\parindent}{1em}
\setlength{\parskip}{0.6em}

\usepackage{graphicx}
\graphicspath{{img/}}          % \includegraphics{filename} looks inside img/
\usepackage{adjustbox}         % lets text auto-scale to fit a given width (e.g. thesis title)

% SVG figures: \includesvg{name} converts img/name.svg to PDF (and PDF_TEX) at
% compile time via inkscape, so editing an SVG only needs a recompile. The
% generated files land in the build directory (see .latexmkrc).
\usepackage{svg}
\svgsetup{
    inkscapeexe=inkscape,
    inkscapearea=drawing,
    inkscapeversion=1,
    inkscapelatex=false,   % render SVG text with inkscape's own fonts, so
                           % underscores etc. in the SVG need no LaTeX escaping
    inkscapepath=build/svg-inkscape/,   % generated PDFs land in the build directory
}
\svgpath{img/}

\usepackage{booktabs}          % \toprule/\midrule/\bottomrule for clean tables
\usepackage{tabularx}          % X columns auto-wrap long content (e.g. file paths)
\usepackage{seqsplit}          % breaks long unbroken strings (paths, tokens) that don't wrap on their own

\usepackage{xcolor}
\newcommand{\todo}[1]{\textbf{\textcolor{red}{[TODO: #1]}}}   % visible marker for unresolved text
\newcommand{\term}[1]{\textsf{#1}}   % hardware/product/protocol names and acronyms (FPGA, DE1-SoC, RGB565, ...)
\newcommand{\sysname}[1]{\textsf{\itshape #1}}   % this thesis's own named systems (FPGAlix, FPGAsteel)

\usepackage[normalem]{ulem}
\newcommand{\reread}[1]{\sout{#1}}   % marks text to go back and review

\usepackage[most]{tcolorbox}
\newtcolorbox{warningbox}{
    colback=red!5!white,
    colframe=red!75!black,
    boxrule=0.6pt,
    arc=2pt,
    left=6pt, right=6pt, top=4pt, bottom=4pt,
}
% Usage: \begin{warningbox} ... \end{warningbox}

\usepackage{hyperref}
\hypersetup{
    colorlinks=true,
    linkcolor=black,
    citecolor=black,
    urlcolor=blue,
    bookmarksnumbered=true,
}

% ------------------------------------------------------------
% C++ code blocks, colors inspired by GitHub's syntax theme
% ------------------------------------------------------------
\usepackage{listings}

\definecolor{githubbg}{HTML}{F6F8FA}
\definecolor{githubkeyword}{HTML}{CF222E}
\definecolor{githubstring}{HTML}{0A3069}
\definecolor{githubcomment}{HTML}{6E7781}
\definecolor{githubnumber}{HTML}{0550AE}
\definecolor{githubrule}{HTML}{D0D7DE}

\lstdefinestyle{githubcpp}{
    language=C++,
    backgroundcolor=\color{githubbg},
    basicstyle=\ttfamily\small,
    keywordstyle=\color{githubkeyword}\bfseries,
    commentstyle=\color{githubcomment}\itshape,
    stringstyle=\color{githubstring},
    numberstyle=\tiny\color{gray},
    identifierstyle=\color{black},
    breaklines=true,
    breakatwhitespace=false,
    columns=fullflexible,
    keepspaces=true,
    showstringspaces=false,
    tabsize=4,
    frame=single,
    rulecolor=\color{githubrule},
    framesep=5pt,
    xleftmargin=4pt,
    xrightmargin=4pt,
    captionpos=b,
    numbers=left,
    numbersep=8pt,
}
\lstset{style=githubcpp}

% Usage in chapters:
%   \begin{lstlisting}[caption={Description}, label={lst:example}]
%   int main() { return 0; }
%   \end{lstlisting}

% ------------------------------------------------------------
% Shell/terminal command blocks, same GitHub-inspired colors
% ------------------------------------------------------------
\lstdefinestyle{githubbash}{
    language=bash,
    backgroundcolor=\color{githubbg},
    basicstyle=\ttfamily\scriptsize,
    keywordstyle=\color{githubkeyword}\bfseries,
    commentstyle=\color{githubcomment}\itshape,
    stringstyle=\color{githubstring},
    identifierstyle=\color{black},
    breaklines=true,
    breakatwhitespace=true,
    columns=fullflexible,
    keepspaces=true,
    showstringspaces=false,
    tabsize=4,
    frame=single,
    rulecolor=\color{githubrule},
    framesep=5pt,
    xleftmargin=4pt,
    xrightmargin=4pt,
    captionpos=b,
    numbers=none,
}
\lstnewenvironment{shellcode}{\lstset{style=githubbash}}{}

% For lines dominated by one long path/token with no useful spaces
% (e.g. export PATH="..."), where breaking only at whitespace causes overflow:
\lstdefinestyle{githubbashpath}{
    language=bash,
    backgroundcolor=\color{githubbg},
    basicstyle=\ttfamily\scriptsize,
    keywordstyle=\color{githubkeyword}\bfseries,
    commentstyle=\color{githubcomment}\itshape,
    stringstyle=\color{githubstring},
    identifierstyle=\color{black},
    breaklines=true,
    breakatwhitespace=false,
    columns=fullflexible,
    keepspaces=true,
    showstringspaces=false,
    tabsize=4,
    frame=single,
    rulecolor=\color{githubrule},
    framesep=5pt,
    xleftmargin=4pt,
    xrightmargin=4pt,
    captionpos=b,
    numbers=none,
}
\lstnewenvironment{shellcodepath}{\lstset{style=githubbashpath}}{}

% Usage in chapters:
%   \begin{shellcode}            <- normal commands / word lists (safe wrapping)
%   ssh wheel@<board-ip>
%   \end{shellcode}
%   \begin{shellcodepath}        <- lines with one long unbroken path/token
%   export PATH="$HOME/very/long/path/bin:$PATH"
%   \end{shellcodepath}

% ------------------------------------------------------------
% Inline code in running text: write |code| between two pipes
% ------------------------------------------------------------
\lstdefinestyle{githubcppinline}{
    language=C++,
    basicstyle=\ttfamily\small\color{black},
    backgroundcolor=\color{githubbg},
    keywordstyle=\color{githubkeyword}\bfseries,
    stringstyle=\color{githubstring},
    showstringspaces=false,
    breaklines=true,
    breakatwhitespace=false,
}
\lstMakeShortInline[style=githubcppinline]|

% Usage in running text:
%   The function |int compute(int x)| returns...
